\chapter{Schwarzschild Metric}

\section*{Introduction}

The Schwarzschild metric maps the abstract geometry of curved spacetime onto measurable physical effects: gravitational time dilation, light bending, orbital precession, and the existence of event horizons. It tells us exactly how mass warps the fabric of spacetime, and it predicts phenomena---like black holes---that were not observed until decades after its derivation.
\section*{Fundamental Equation Used}

\begin{equation}
ds^2 = -\left(1 - \frac{2GM}{rc^2}\right)c^2\,dt^2 + \left(1 - \frac{2GM}{rc^2}\right)^{-1}dr^2 + r^2\,d\Omega^2
\end{equation}
\section*{Assumptions}

\textbf{Assumptions:} Spherical symmetry; vacuum outside the mass (no charge, no rotation); static (time-independent) metric.


\textbf{Limitations:} Does not describe the interior of the mass; rotation (Kerr metric) and charge (Reissner-Nordstrom metric) require generalizations.


\section*{Mathematical Derivation}

\textbf{Step 1: Assume a static, spherically symmetric metric.}

By Birkhoff's theorem, any spherically symmetric vacuum solution of Einstein's field equations must be static. The most general static, spherically symmetric metric in Schwarzschild coordinates $(t, r, \theta, \phi)$ is
\begin{equation}
ds^2 = -A(r)\,c^2\,dt^2 + B(r)\,dr^2 + r^2\,d\Omega^2,
\end{equation}
where $A(r)$ and $B(r)$ are unknown functions and $d\Omega^2 = d\theta^2 + \sin^2\theta\,d\phi^2$.

\textbf{Step 2: Impose the vacuum Einstein equations $R_{\mu\nu} = 0$.}

Computing the Ricci tensor components for this metric and setting them to zero yields the system:
\begin{align}
R_{tt} &= \frac{A''}{2B} - \frac{A'}{4B}\left(\frac{A'}{A} + \frac{B'}{B}\right) + \frac{A'}{rB} = 0, \\
R_{rr} &= -\frac{A''}{2A} + \frac{A'}{4A}\left(\frac{A'}{A} + \frac{B'}{B}\right) + \frac{B'}{rB} = 0, \\
R_{\theta\theta} &= 1 - \frac{1}{B} + \frac{r}{2B}\left(\frac{A'}{A} - \frac{B'}{B}\right) = 0.
\end{align}

\textbf{Step 3: Solve the system of ODEs.}

From $R_{tt}/A + R_{rr}/B = 0$, one obtains $(A'B + AB')/(AB) = A'/A - B'/B$, which simplifies to $A'B/AB + B'/B = A'/A$. This implies $A' B' = 0$ in combination with the other equations. A cleaner approach: adding $R_{tt}/A$ and $R_{rr}/B$ gives $(A'B + AB')/(rAB) = 0$, hence $A(r)B(r) = \text{const}$. Asymptotic flatness ($r\to\infty$: $A\to 1$, $B\to 1$) fixes the constant to 1:
\begin{equation}
B(r) = \frac{1}{A(r)}.
\end{equation}

\textbf{Step 4: Substitute and solve for $A(r)$.}

Substituting $B = 1/A$ into $R_{\theta\theta} = 0$:
\begin{equation}
1 - A + \frac{r}{2}\left(A' \cdot \frac{1}{A}\cdot A + 0\right) \cdot \frac{1}{A} \;\longrightarrow\; 1 - A + rA' = 0.
\end{equation}
(Working more carefully with the substitution $B = 1/A$ and the original $R_{\theta\theta}$ equation.) This ODE has the solution:
\begin{equation}
A(r) = 1 - \frac{2GM}{c^2\,r},
\end{equation}
where $2GM/c^2$ is the integration constant, fixed by matching the Newtonian limit $g_{tt} \to -(1 + 2\Phi/c^2)$ with $\Phi = -GM/r$.

\textbf{Step 5: Assemble the Schwarzschild metric.}

With $A(r) = B(r)^{-1} = 1 - 2GM/(c^2 r)$, the full metric is
\begin{equation}
\boxed{ds^2 = -\left(1 - \frac{2GM}{rc^2}\right)c^2\,dt^2 + \left(1 - \frac{2GM}{rc^2}\right)^{-1}dr^2 + r^2\,d\Omega^2.}
\end{equation}
The Schwarzschild radius $r_s = 2GM/c^2$ is the event horizon; the true singularity lies at $r = 0$.

\section*{Characteristics of the Equation}

\begin{itemize}
    \item The metric has a coordinate singularity at $r = r_s$, the event horizon.
    \item $d\Omega^2 = d\theta^2 + \sin^2\theta\,d\phi^2$ is the metric on the unit two-sphere.
    \item For $r \gg r_s$, the metric reduces to the flat Minkowski metric (weak-field limit).
    \item Derived by Karl Schwarzschild in 1916, shortly after Einstein published general relativity.
\end{itemize}
\section*{Solution Methods}

\begin{itemize}
    \item \textbf{Direct substitution:} Verify the metric satisfies $R_{\mu\nu} = 0$ (vacuum field equations).
    \item \textbf{Geodesic calculation:} Derive orbital equations from the effective potential $V_{eff} = -\frac{GM}{r} + \frac{L^2}{2r^2} - \frac{GML^2}{c^2r^3}$.
    \item \textbf{Redshift factor:} Compute gravitational time dilation using $d\tau = \sqrt{1 - r_s/r}\,dt$.
\end{itemize}
\section*{Verifications}

\begin{itemize}
    \item Mercury's perihelion precession of 43 arcseconds per century is correctly predicted.
    \item Light deflection by the Sun (1.75 arcseconds) matches the 1919 Eddington observations.
    \item Gravitational time dilation has been measured with atomic clocks at different altitudes.
\end{itemize}
\section*{Applications}

\begin{itemize}
    \item Gravitational redshift predictions confirmed by Pound-Rebka experiment (1959).
    \item Black hole event horizons and singularities.
    \item Gravitational lensing and light deflection by massive objects.
    \item Orbital precession of Mercury and binary pulsars.
\end{itemize}
\section*{What Richard Feynman Would Have Asked}

\subsection*{Plain-English Translation}
\textit{``If you had to explain the Schwarzschild metric to someone who has never studied general relativity, what would you say?''}

The Schwarzschild metric is a precise mathematical description of how a spherical mass like a star warps the space and time around it. Time runs slower near the mass, radial distances are stretched, and light follows curved paths. The metric is the exact solution to Einstein's gravity equations, and it tells you everything about the geometry of spacetime outside a star---including the existence of black holes, where the warping becomes so extreme that nothing can escape.

\subsection*{Localized Mechanics}
\textit{``What is actually happening to space and time at a specific point near a massive object? How does the metric describe local physics?''}

At any specific point, the metric tells you two things: how fast your clock ticks relative to a distant observer (the $g_{tt}$ component), and how a ruler measures radial distance (the $g_{rr}$ component). Near a massive object, clocks tick slower and radial distances are stretched. A local observer with a small enough laboratory cannot tell the difference from flat spacetime (the equivalence principle), but a distant observer sees these effects clearly. The metric encodes these local distortions as functions of $r$.

\subsection*{Testing Extreme Limits}
\textit{``What happens in the extreme limit as you approach the singularity at $r = 0$? Does the metric really predict infinite curvature?''}

Yes, at $r = 0$ the metric predicts infinite curvature---a true physical singularity where general relativity breaks down. The tidal forces become infinite, and the classical description of spacetime ceases to be meaningful. This is widely believed to require a quantum theory of gravity for resolution. Near the singularity, quantum effects (which are not included in the classical Schwarzschild metric) are expected to prevent or modify the singular behavior.

\subsection*{Dimensionality and Geometry}
\textit{``How would the Schwarzschild solution change if the universe had more than three spatial dimensions?''}

In $d$ spatial dimensions, the gravitational potential falls off as $1/r^{d-2}$ instead of $1/r$, and the Schwarzschild metric generalizes accordingly. In $d = 4$ (four spatial dimensions), the event horizon structure changes, and the singularity at $r = 0$ becomes ``softer'' (the tidal forces diverge more slowly). The existence of extra dimensions would modify black hole thermodynamics, gravitational wave emission, and the structure of singularities---all active areas of research in string theory and brane-world scenarios.

\subsection*{Cross-Disciplinary Connection}
\textit{``Where does the same mathematical structure---a metric with a coordinate singularity and a true singularity---appear in other fields?''}

The Schwarzschild metric's mathematical structure is echoed in the ergosphere of rotating black holes (Kerr metric), in the sonic analog of black holes (``dumb holes'') in fluid dynamics, and in optical systems where light cannot escape (optical black holes). The coordinate singularity at the event horizon is analogous to a coordinate transformation that becomes degenerate, like the transformation from Cartesian to polar coordinates at the origin. The deep mathematical lesson is that coordinate singularities are artifacts of our description, while true singularities reflect fundamental physical limits.

% ------------------------------------------------------------
\section*{FAQ}

\textbf{Q: What happens at the Schwarzschild radius?}
A: The Schwarzschild radius $r_s = 2GM/c^2$ is the event horizon. For a distant observer, nothing crosses it---signals are infinitely redshifted. For an infalling observer, nothing special happens locally at $r_s$; the observer passes through smoothly. The ``singularity'' at $r_s$ is a coordinate artifact, not a physical one.

\textbf{Q: Can anything escape from inside the event horizon?}
A: No. Inside $r_s$, the roles of space and time effectively swap: the radial direction becomes timelike, and moving toward $r = 0$ is as inevitable as moving forward in time. No signal, not even light, can escape to $r > r_s$.

\textbf{Q: Does the Schwarzschild metric apply to the Sun?}
A: Approximately yes. The Sun is nearly spherical and slowly rotating, so the Schwarzschild metric is a good approximation outside the Sun. For precise calculations, small corrections from the Sun's rotation (frame-dragging) are needed.
\section*{Important Bibliography}

\begin{enumerate}
\item Misner, C.W., Thorne, K.S. and Wheeler, J.A., \textit{Gravitation} (W.H.\ Freeman, 1973), Chapters 18 and 23.
\item Schutz, B.F., \textit{A First Course in General Relativity}, 2nd ed. (Cambridge University Press, 2009), Chapter 11.
\item Carroll, S.M., \textit{Spacetime and Geometry: An Introduction to General Relativity} (Cambridge University Press, 2019), Chapter 5.
\item Weinberg, S., \textit{Gravitation and Cosmology} (Wiley, 1972), Chapter 8.
\end{enumerate}


